\section{Why Volatility Is Not a Sufficient State}

\subsection{Proxy correlation is not decision stability}

The two monthly volatility series have a correlation of 0.83. If volatility were used as
a continuous descriptive variable, this might appear reassuring. The router, however,
uses a discontinuous decision: crossing the rolling 75th percentile changes the entire
portfolio from momentum to reversal. Small differences in weighting or rank near the
boundary can therefore produce different holdings.

Across the 66 holding months, the two proxies assign the same quartile only 42.4\% of the
time. Their Q4 sets have a Jaccard overlap of 54.5\%. \Cref{fig:stability} distinguishes
correlation in levels from agreement in the actual decision state.

\begin{figure}[H]
\centering
\includegraphics[width=0.78\textwidth]{fig03_public_proxy_stability.png}
\caption{Aggregate agreement between the two pre-specified futures-market volatility
proxies.}
\label{fig:stability}
\end{figure}

This disagreement is not a technical nuisance. Open-interest-notional weighting asks
where outstanding notional activity is concentrated; equal weighting asks what happens
to the typical product. Futures have no theorem that makes either object the unique
``market portfolio.'' A successful router must either survive both or provide an
economic reason why one proxy is causally tied to sleeve failure.

\subsection{High volatility does identify stronger common shocks}

The failure is not explained by Q4 being random. Under the primary proxy, average
pairwise product correlation rises from 0.147 in Q1 to 0.254 in Q4. Average
cross-sectional dispersion rises from 1.21\% to 1.71\%, and directional coherence rises
from 0.300 to 0.376. Q4 therefore represents a more systemic and more heterogeneous
shock environment.

\begin{figure}[H]
\centering
\includegraphics[width=0.96\textwidth]{fig04_public_market_anatomy.png}
\caption{Cross-sectional market anatomy in the primary proxy's signal month.}
\label{fig:anatomy}
\end{figure}

The simultaneous increase in correlation and dispersion is revealing. More products move
in a common direction, while the magnitude of their responses becomes more unequal. This
can arise when a common macro shock is transmitted through products with different
exposures, or when several large sector shocks coincide. Neither case determines whether
the cross-sectional ranking will reverse next month.

\subsection{Intensity is a scalar; strategy state is a vector}

Volatility answers how large price changes are. A momentum--reversal decision additionally
depends on at least four dimensions:

\begin{enumerate}
  \item \textbf{Direction:} a distressed decline followed by rebound versus a directional
  repricing that continues;
  \item \textbf{Source:} demand, supply, inventory, policy, funding, currency, or liquidity;
  \item \textbf{Breadth:} broad participation versus concentration in a subset of large
  products;
  \item \textbf{Stage:} shock initiation, continued absorption, crowded liquidation, or
  exhaustion and repair.
\end{enumerate}

Two months can have the same $\sigma_t$ and differ along every dimension above. The
equity result can be understood as a setting in which high volatility more reliably
coincides with a particular distress-and-rebound vector. The futures cross-section mixes
trend-generating physical shocks with reversal-generating liquidity shocks. A scalar
cannot identify which one is present.
